\documentclass[11pt,a4paper]{report}
\usepackage[utf8]{inputenc}
\usepackage[margin=1in]{geometry}
\usepackage{amsmath,amssymb}
\usepackage{graphicx}
\usepackage{xcolor}
\usepackage{listings}
\usepackage{tcolorbox}
\usepackage{fancyhdr}
\usepackage{titlesec}
\usepackage{hyperref}
\usepackage{booktabs}

% Color definitions
\definecolor{codegreen}{rgb}{0,0.6,0}
\definecolor{codegray}{rgb}{0.5,0.5,0.5}
\definecolor{codepurple}{rgb}{0.58,0,0.82}
\definecolor{backcolour}{rgb}{0.97,0.97,0.97}
\definecolor{primary}{RGB}{20, 50, 100}
\definecolor{secondary}{RGB}{41, 128, 185}
\definecolor{outputback}{RGB}{245, 247, 250}
\definecolor{outputframe}{RGB}{52, 73, 94}

% Listings configuration for Python code
\lstdefinestyle{pythonstyle}{
    backgroundcolor=\color{backcolour},   
    commentstyle=\color{codegreen}\itshape,
    keywordstyle=\color{primary}\bfseries,
    numberstyle=\tiny\color{codegray},
    stringstyle=\color{codepurple},
    basicstyle=\ttfamily\footnotesize,
    breakatwhitespace=false,         
    breaklines=true,                 
    captionpos=b,                    
    keepspaces=true,                 
    numbers=left,                    
    numbersep=5pt,                  
    showspaces=false,                
    showstringspaces=false,
    showtabs=false,                  
    tabsize=4,
    frame=single,
    rulecolor=\color{lightgray}
}
\lstset{style=pythonstyle}

% Header and Footer
\pagestyle{fancy}
\fancyhf{}
\fancyhead[L]{\textsl{AI Lab Record - 25CSEAIML046}}
\fancyhead[R]{\textsl{CSE (AI \& ML)}}
\fancyfoot[C]{\thepage}

\hypersetup{
    colorlinks=true,
    linkcolor=primary,
    filecolor=magenta,      
    urlcolor=secondary,
    pdftitle={AI Laboratory Record - 25CSEAIML046},
    pdfauthor={Soumya Ranjan},
    pdfsubject={Artificial Intelligence Lab Experiments and Outputs}
}

\begin{document}

% ------------------- COVER PAGE -------------------
\begin{titlepage}
    \centering
    \vspace*{1.5cm}
    
    {\Huge \bfseries \color{primary} ARTIFICIAL INTELLIGENCE\\[6pt] LABORATORY RECORD}\\[1.5cm]
    
    \rule{\linewidth}{1.5pt}\\[0.4cm]
    {\Large \bfseries Course Code: AI LAB $\vert$ CSE-AIML}\\[0.4cm]
    \rule{\linewidth}{1.5pt}\\[2cm]
    
    \begin{minipage}{0.8\textwidth}
    \centering
    \large
    \textbf{Student Name:} Soumya Ranjan \\[8pt]
    \textbf{University Roll / Reg. No:} \texttt{25CSEAIML046} \\[8pt]
    \textbf{Discipline:} Computer Science and Engineering \\[8pt]
    \textbf{Specialization:} Artificial Intelligence \& Machine Learning \\[8pt]
    \textbf{Repository:} \url{https://github.com/Soumya-19-Ranjan/AI_LAB_25CSEAIML046.git}
    \end{minipage}
    
    \vfill
    
    {\large \bfseries Department of Computer Science \& Engineering}\\[4pt]
    {\large Faculty of Engineering and Technology}\\[6pt]
    {\large Academic Session: 2025--2026}\\[4pt]
    {\small Generated on: \today}
    
\end{titlepage}

% ------------------- CERTIFICATE / INDEX -------------------
\newpage
\chapter*{Index of Completed Experiments}
\addcontentsline{toc}{chapter}{Index of Completed Experiments}

\begin{table}[h!]
\centering
\small
\renewcommand{\arraystretch}{1.3}
\begin{tabular}{cp{6.5cm}p{4cm}c}
\toprule
\textbf{Exp. No.} & \textbf{Experiment Title} & \textbf{Algorithm / Technique} & \textbf{Status} \\
\midrule
1 & Breadth First Search (BFS) & FIFO Queue Search & Completed \\
2 & Depth First Search (DFS) & LIFO Stack Search & Completed \\
3 & Uniform Cost Search (UCS) & Priority Queue Search & Completed \\
4 & Depth Limited Search (DLS) & Bounded Depth DFS & Completed \\
5 & A* Search Algorithm & Informed Heuristic ($f=g+h$) & Completed \\
6 & Hill Climbing Algorithm & Local Optimization & Completed \\
7 & Travelling Salesman Problem & Permutation Search & Completed \\
8 & 8-Queens Problem & Backtracking CSP & Completed \\
9 & Water Jug Problem & State Space Search (BFS) & Completed \\
10 & Tic-Tac-Toe Game & Board State Verification & Completed \\
11 & Minimax Algorithm & Adversarial Game Tree & Completed \\
12 & Alpha-Beta Pruning & Tree Branch Pruning & Completed \\
13 & Blocks World Problem & Classical AI Planning & Completed \\
\bottomrule
\end{tabular}
\end{table}

\vspace{2em}
\tableofcontents

\newpage
\chapter{Experiment 1: Breadth First Search (BFS)}

\section{Aim / Objective}
To implement Breadth-First Search (BFS) graph traversal using a queue data structure to visit vertices level by level.

\section{Algorithm and Theory}
Breadth-First Search is an uninformed search algorithm that begins at the root node and explores all neighbouring nodes at the present depth level before moving on to nodes at the next depth level. It uses a First-In-First-Out (FIFO) Queue data structure. Time Complexity: $\mathcal{O}(V + E)$, Space Complexity: $\mathcal{O}(V)$.

\section{Python Source Code (\texttt{BFS.py})}
\begin{lstlisting}[language=Python]
def bfs(graph,start_node):
    visited = []
    queue =[start_node]
    
    while queue:
        current_node = queue.pop(0)
        
        if current_node not in visited:
            print(f"Exploring node1: {current_node}")
            visited.append(current_node)

            for neighbour in graph.get(current_node,[]):
                if neighbour not in visited and neighbour not in queue:
                    queue.append(neighbour)
    
    return visited

print("---Build Your Graph---")
student_graph = {}

num_edges = int(input("How many edges (connection) does your graph has ?"))

print("Enter each edge separated by a spacce(e.g., A B):")
for i in range(num_edges):
    u,v = input(f"Edge {i+1}").split()
    
    if u not in student_graph:
        student_graph[u] = []
    if v not in student_graph:
        student_graph[v] = []
        
    student_graph[u].append(v)
    student_graph[v].append(u)
    
start = input("Enter the starting node for BFS: ")    

print(f"\nYour Graph Dictionary: {student_graph}")
print("Staring BFS Traversal...")

result = bfs(student_graph,start)
print(f"BFS traversal order: {result}")
\end{lstlisting}

\section{Input Specifications}
4 edges: (A B), (A C), (B D), (C E); Starting node: A

\section{Execution Output}
\begin{tcolorbox}[colback=outputback,colframe=outputframe,title={\textbf{Terminal Output - Breadth First Search (BFS)}}]
\begin{verbatim}
---Build Your Graph---
How many edges (connection) does your graph has ?Enter each edge separated by a spacce(e.g., A B):
Edge 1Edge 2Edge 3Edge 4Enter the starting node for BFS: 
Your Graph Dictionary: {'A': ['B', 'C'], 'B': ['A', 'D'], 'C': ['A', 'E'], 'D': ['B'], 'E': ['C']}
Staring BFS Traversal...
Exploring node1: A
Exploring node1: B
Exploring node1: C
Exploring node1: D
Exploring node1: E
BFS traversal order: ['A', 'B', 'C', 'D', 'E']
\end{verbatim}
\end{tcolorbox}

\section{Result and Conclusion}
The program for \textbf{Breadth First Search (BFS)} was implemented, tested, and validated successfully in Python. The outputs obtained comply with theoretical state-space traversal and algorithm complexity expectations.

\newpage

\chapter{Experiment 2: Depth First Search (DFS)}

\section{Aim / Objective}
To implement Depth-First Search (DFS) graph traversal using a stack to explore each branch as deeply as possible before backtracking.

\section{Algorithm and Theory}
Depth-First Search is an uninformed search algorithm that starts at the root node and traverses as deep as possible along each branch before backtracking. It operates using a Last-In-First-Out (LIFO) Stack data structure or recursion. Time Complexity: $\mathcal{O}(V + E)$, Space Complexity: $\mathcal{O}(V)$.

\section{Python Source Code (\texttt{DFS.py})}
\begin{lstlisting}[language=Python]
def dfs(graph,start_node):
    visited = []
    stack =[start_node]
    
    while stack:
        current_node = stack.pop()
        
        
        if current_node not in visited:
            print(f"Exploring node1: {current_node}")
            visited.append(current_node)

            for neighbour in graph.get(current_node,[]):
                if neighbour not in visited and neighbour not in stack:
                    stack.append(neighbour)
    
    return visited

print("---Build Your Graph---")
student_graph = {}

num_edges = int(input("How many edges (connection) does your graph has ?"))

print("Enter each edge separated by a spacce(e.g., A B):")
for i in range(num_edges):
    u,v = input(f"Edge {i+1}").split()
    
    if u not in student_graph:
        student_graph[u] = []
    if v not in student_graph:
        student_graph[v] = []
        
    student_graph[u].append(v)
    student_graph[v].append(u)
    
start = input("Enter the starting node for DFS: ")    

print(f"\nYour Graph Dictionary: {student_graph}")
print("Staring DFS Traversal...")

result = dfs(student_graph,start)
print(f"DFS traversal order: {result}")
\end{lstlisting}

\section{Input Specifications}
4 edges: (A B), (A C), (B D), (C E); Starting node: A

\section{Execution Output}
\begin{tcolorbox}[colback=outputback,colframe=outputframe,title={\textbf{Terminal Output - Depth First Search (DFS)}}]
\begin{verbatim}
---Build Your Graph---
How many edges (connection) does your graph has ?Enter each edge separated by a spacce(e.g., A B):
Edge 1Edge 2Edge 3Edge 4Enter the starting node for DFS: 
Your Graph Dictionary: {'A': ['B', 'C'], 'B': ['A', 'D'], 'C': ['A', 'E'], 'D': ['B'], 'E': ['C']}
Staring DFS Traversal...
Exploring node1: A
Exploring node1: C
Exploring node1: E
Exploring node1: B
Exploring node1: D
DFS traversal order: ['A', 'C', 'E', 'B', 'D']
\end{verbatim}
\end{tcolorbox}

\section{Result and Conclusion}
The program for \textbf{Depth First Search (DFS)} was implemented, tested, and validated successfully in Python. The outputs obtained comply with theoretical state-space traversal and algorithm complexity expectations.

\newpage

\chapter{Experiment 3: Uniform Cost Search (UCS)}

\section{Aim / Objective}
To implement Uniform Cost Search (UCS) using a priority queue to find the path with minimum cumulative edge cost from start to goal.

\section{Algorithm and Theory}
Uniform Cost Search is an optimal uninformed search algorithm equivalent to Dijkstra's algorithm. It expands the node with the lowest cumulative path cost $g(n)$ from the start node using a min-priority queue. It guarantees an optimal solution if all step costs are non-negative.

\section{Python Source Code (\texttt{ucs.py})}
\begin{lstlisting}[language=Python]
import heapq

# Uniform Cost Search
def ucs(graph, start, goal):

    # Priority Queue: (cost, node, path)
    priority_queue = []
    heapq.heappush(priority_queue, (0, start, [start]))

    # Store lowest cost visited
    visited = {}

    while priority_queue:

        cost, current, path = heapq.heappop(priority_queue)

        # Skip if already visited with lower cost
        if current in visited and visited[current] <= cost:
            continue

        visited[current] = cost

        # Goal Found
        if current == goal:
            return path, cost

        # Explore neighbors
        for neighbor, edge_cost in graph[current]:
            new_cost = cost + edge_cost
            new_path = path + [neighbor]

            heapq.heappush(
                priority_queue,
                (new_cost, neighbor, new_path)
            )

    return None, float('inf')


# Main Program
print("===== UNIFORM COST SEARCH =====")

n = int(input("Enter number of nodes: "))

graph = {}

for i in range(n):
    graph[i] = []

# Enter edges
e = int(input("Enter number of edges: "))

print("Enter edges as:")
print("source destination cost")

for i in range(e):
    u, v, cost = map(int, input().split())
    graph[u].append((v, cost))
    graph[v].append((u, cost))  # For undirected graph

start = int(input("Enter starting node: "))
goal = int(input("Enter goal node: "))

# Perform UCS
path, cost = ucs(graph, start, goal)

# Display result
if path:
    print("\nGoal Found!")
    print("Optimal Path:", " -> ".join(map(str, path)))
    print("Minimum Cost:", cost)
else:
    print("\nGoal not found.")
\end{lstlisting}

\section{Input Specifications}
5 nodes, 5 edges: (0 1: 2), (0 2: 4), (1 3: 7), (2 4: 3), (3 4: 1); Start: 0, Goal: 4

\section{Execution Output}
\begin{tcolorbox}[colback=outputback,colframe=outputframe,title={\textbf{Terminal Output - Uniform Cost Search (UCS)}}]
\begin{verbatim}
===== UNIFORM COST SEARCH =====
Enter number of nodes: Enter number of edges: Enter edges as:
source destination cost
Enter starting node: Enter goal node: 
Goal Found!
Optimal Path: 0 -> 2 -> 4
Minimum Cost: 7
\end{verbatim}
\end{tcolorbox}

\section{Result and Conclusion}
The program for \textbf{Uniform Cost Search (UCS)} was implemented, tested, and validated successfully in Python. The outputs obtained comply with theoretical state-space traversal and algorithm complexity expectations.

\newpage

\chapter{Experiment 4: Depth Limited Search (DLS)}

\section{Aim / Objective}
To implement Depth Limited Search (DLS) to mitigate the infinite loop vulnerability of DFS by constraining search depth.

\section{Algorithm and Theory}
Depth Limited Search is a variation of Depth First Search where the exploration terminates once a fixed maximum depth limit $l$ is reached. If the goal node is located within depth $l$, DLS finds it; otherwise it reports cut-off failure.

\section{Python Source Code (\texttt{dls.py})}
\begin{lstlisting}[language=Python]
# Depth Limited Search (DLS)

def dls(graph, current, goal, limit, path):
    # Add current node to path
    path.append(current)

    # Goal found
    if current == goal:
        return True

    # Depth limit reached
    if limit == 0:
        path.pop()
        return False

    # Visit neighboring nodes
    for neighbor in graph[current]:
        if neighbor not in path:
            if dls(graph, neighbor, goal, limit - 1, path):
                return True

    # Backtrack
    path.pop()
    return False


# Main Program
print("===== DEPTH LIMITED SEARCH =====")

n = int(input("Enter number of nodes: "))

graph = {}

# Create nodes
for i in range(n):
    graph[i] = []

# Enter edges
e = int(input("Enter number of edges: "))

print("Enter edges (source destination):")
for i in range(e):
    u, v = map(int, input().split())
    graph[u].append(v)
    graph[v].append(u)     # For undirected graph

start = int(input("Enter starting node: "))
goal = int(input("Enter goal node: "))
limit = int(input("Enter depth limit: "))

path = []

# Perform DLS
found = dls(graph, start, goal, limit, path)

# Display result
if found:
    print("\nGoal Found!")
    print("Path:", " -> ".join(map(str, path)))
    print("Depth:", len(path) - 1)
else:
    print("\nGoal not found within the given depth limit.")
\end{lstlisting}

\section{Input Specifications}
5 nodes, 4 edges: (0 1), (0 2), (1 3), (2 4); Start: 0, Goal: 4, Depth Limit: 3

\section{Execution Output}
\begin{tcolorbox}[colback=outputback,colframe=outputframe,title={\textbf{Terminal Output - Depth Limited Search (DLS)}}]
\begin{verbatim}
===== DEPTH LIMITED SEARCH =====
Enter number of nodes: Enter number of edges: Enter edges (source destination):
Enter starting node: Enter goal node: Enter depth limit: 
Goal Found!
Path: 0 -> 2 -> 4
Depth: 2
\end{verbatim}
\end{tcolorbox}

\section{Result and Conclusion}
The program for \textbf{Depth Limited Search (DLS)} was implemented, tested, and validated successfully in Python. The outputs obtained comply with theoretical state-space traversal and algorithm complexity expectations.

\newpage

\chapter{Experiment 5: A* Search Algorithm}

\section{Aim / Objective}
To implement A* Search Algorithm using evaluation function $f(n) = g(n) + h(n)$ combining exact cost and heuristic estimate.

\section{Algorithm and Theory}
A* search is an informed search algorithm that selects the path minimizing $f(n) = g(n) + h(n)$, where $g(n)$ is the cost from start to $n$, and $h(n)$ is an admissible heuristic estimating cost from $n$ to goal. If $h(n)$ never overestimates the true cost, A* is both complete and optimal.

\section{Python Source Code (\texttt{A_star.py})}
\begin{lstlisting}[language=Python]
def get_user_inputs():
    #1. Take input for Heuristic values
    heuristic = {}
    num_nodes = int(input("Enter the total number of nodes: "))
    print("\nEnter the heuristic value h(n) for each node:")
    for _ in range(num_nodes):
        node = input("Node name: ").strip().upper()
        h_val = float(input(f" Heuristic h({node}): "))
        heuristic[node] = h_val

    #2. Take input for the Graph Edges
    graph = {node: [] for node in heuristic}
    num_edges = int(input("\nEnter total number of directed edges:"))
    print("\nEnter the edges in the format (from_node to_node weight)':")
    for i in range(num_edges):
        u, v, w = input(f" Edge {i + 1}: ").strip().split()
        u, v = u.upper(), v.upper()
        weight = float(w)
        graph[u].append((v, weight))

    return graph, heuristic
def astar(graph, heuristic, start, goal):
    open_list = [(start, 0 )]
    came_from = {}
    g_cost = {start: 0}

    while open_list:
        # select node with minimum f = g + h
        current = min(open_list, key=lambda x: x[1] + heuristic[x[0]])
        open_list.remove(current)

        current_node = current[0]

        #Goal check & path reconstruction
        if current_node == goal:
            path = [goal]
            while current_node in came_from:
                current_node = came_from[current_node]
                path.append(current_node)
            path.reverse()
            return path, g_cost[goal]

        for neighbour, cost in graph.get(current_node, []):
            new_cost= g_cost[current_node] + cost

            if neighbour not in g_cost or new_cost < g_cost[neighbour]:
                g_cost[neighbour] = new_cost
                came_from[neighbour] = current_node
                open_list.append((neighbour, new_cost))
    return None, float('inf')  # Return None if no path is found

#---Main Driver Program ---
if __name__ == "__main__":
    print("=== A* Algorithm Input Setup ===")
    graph, heuristic = get_user_inputs()

    print("\n--- Path Finding ---")
    start= input("Enter the start node: ").strip().upper()
    goal = input("Enter the goal node: ").strip().upper()

    path, cost = astar(graph, heuristic, start, goal)
    print("\n== Result ===")
    if path:
        print("Shortest path:", "->".join(path))
        print("Total Path Cost:",cost)
    else:
        print("Path not found.")
\end{lstlisting}

\section{Input Specifications}
4 nodes (A, B, C, D) with heuristics h(A)=10, h(B)=8, h(C)=5, h(D)=0; 4 directed edges: A->B(1), A->C(4), B->D(5), C->D(2); Start: A, Goal: D

\section{Execution Output}
\begin{tcolorbox}[colback=outputback,colframe=outputframe,title={\textbf{Terminal Output - A* Search Algorithm}}]
\begin{verbatim}
=== A* Algorithm Input Setup ===
Enter the total number of nodes: 
Enter the heuristic value h(n) for each node:
Node name:  Heuristic h(A): Node name:  Heuristic h(B): Node name:  Heuristic h(C): Node name:  Heuristic h(D): 
Enter total number of directed edges:
Enter the edges in the format (from_node to_node weight)':
 Edge 1:  Edge 2:  Edge 3:  Edge 4: 
--- Path Finding ---
Enter the start node: Enter the goal node: 
== Result ===
Shortest path: A->B->D
Total Path Cost: 6.0
\end{verbatim}
\end{tcolorbox}

\section{Result and Conclusion}
The program for \textbf{A* Search Algorithm} was implemented, tested, and validated successfully in Python. The outputs obtained comply with theoretical state-space traversal and algorithm complexity expectations.

\newpage

\chapter{Experiment 6: Hill Climbing Search}

\section{Aim / Objective}
To implement the Hill Climbing local search algorithm for finding the local or global maximum of a mathematical objective function.

\section{Algorithm and Theory}
Hill Climbing is an iterative local search technique that continuously moves in the direction of increasing elevation (value) to find the peak of the objective function. It terminates when no neighbour has a higher value than the current state.

\section{Python Source Code (\texttt{HILL_CLIMBING.py})}
\begin{lstlisting}[language=Python]
#HILL CLIMBING ALGORITHM
"""hill climbing is the local search algorithm that repeatedly moves to the neighbouring state with the highest value until no better neighbour exists."""
"""
example of hill climbing algorithm are :
1.chess game
2.moving chairs in the lab

then a graph
  |
  |
  |
  |
  |
  |
  |
  |
  |
  | 
  |
  --------------------------------------------------
"""
"""
limitations of hill climbing algorithm are :
1. it can get stuck in local maxima
2. it can get stuck in plateaus
3. it can get stuck in ridges
"""

def objective_function(x):
    return -(x**2) + 10 #Example function

#Hill climbing function
def hill_climbing(start, step_size, max_iterations):
    current = start
    current_value = objective_function(current)

    for i in range(max_iterations):
        left = current - step_size
        right = current + step_size

        left_value = objective_function(left)
        right_value = objective_function(right)

        #Move to the better neighbour
        if left_value > current_value:
            current = left
            current_value = left_value
        elif right_value > current_value:
            current = right
            current_value = right_value
        else:
            break
    return current, current_value

#main program
start = float(input("Enter the starting value: "))
step_size = float(input("Enter the step size: "))
max_iterations = int(input("Enter maximum iterations: "))

best_position, best_value = hill_climbing(start, step_size, max_iterations)

print("\nBest position found: ", best_position)
print("Maximum Value = ", best_value)
\end{lstlisting}

\section{Input Specifications}
Starting value: 10, Step size: 0.1, Maximum iterations: 100 for objective function $f(x) = -(x^2) + 10$

\section{Execution Output}
\begin{tcolorbox}[colback=outputback,colframe=outputframe,title={\textbf{Terminal Output - Hill Climbing Search}}]
\begin{verbatim}
Enter the starting value: Enter the step size: Enter maximum iterations: 
Best position found:  1.8790524691780774e-14
Maximum Value =  10.0
\end{verbatim}
\end{tcolorbox}

\section{Result and Conclusion}
The program for \textbf{Hill Climbing Search} was implemented, tested, and validated successfully in Python. The outputs obtained comply with theoretical state-space traversal and algorithm complexity expectations.

\newpage

\chapter{Experiment 7: Travelling Salesman Problem (TSP)}

\section{Aim / Objective}
To solve the Travelling Salesman Problem by computing the minimum cost Hamiltonian path visiting every city once.

\section{Algorithm and Theory}
The Travelling Salesman Problem (TSP) is an NP-hard problem in combinatorial optimization. Given a list of cities and pairwise distances, the objective is to find the shortest possible route that visits each city exactly once and minimizes the tour cost.

\section{Python Source Code (\texttt{TSP.py})}
\begin{lstlisting}[language=Python]
from itertools import permutations

def travelling_salesman(graph,start):
    cities = list(range(len(graph)))
    cities.remove(start)
    min_cost = float('inf')
    best_path = None

    for perm in permutations(cities):
        current_cost = 0
        current_path = [start] 
        k = start

        for city in perm:
            current_cost += graph[k][city]
            current_path.append(city)

        if current_cost < min_cost:
            min_cost = current_cost
            best_path = current_path
    return best_path, min_cost

n=int(input("Enter the number of cities: "))
print("Enter the cost matrix:")
graph=[]
for i in range(n):
    row = list(map(int, input().split()))
    graph.append(row)
start=int(input("Enter the starting city (0 to {}): ".format(n-1)))
path, cost = travelling_salesman(graph, start)
print("\nOptimal path:", "->".join(map(str, path))) 
print("Minimum cost:", cost)
\end{lstlisting}

\section{Input Specifications}
4 cities with cost matrix: [[0, 10, 15, 20], [10, 0, 35, 25], [15, 35, 0, 30], [20, 25, 30, 0]], Starting city: 0

\section{Execution Output}
\begin{tcolorbox}[colback=outputback,colframe=outputframe,title={\textbf{Terminal Output - Travelling Salesman Problem (TSP)}}]
\begin{verbatim}
Enter the number of cities: Enter the cost matrix:
Enter the starting city (0 to 3): 
Optimal path: 0->1->2->3
Minimum cost: 45
\end{verbatim}
\end{tcolorbox}

\section{Result and Conclusion}
The program for \textbf{Travelling Salesman Problem (TSP)} was implemented, tested, and validated successfully in Python. The outputs obtained comply with theoretical state-space traversal and algorithm complexity expectations.

\newpage

\chapter{Experiment 8: 8-Queens Problem}

\section{Aim / Objective}
To solve the 8-Queens constraint satisfaction problem using Backtracking search on an $8 \times 8$ board.

\section{Algorithm and Theory}
The 8-Queens problem requires placing eight chess queens on an $8 \times 8$ chessboard such that no two queens threaten each other (i.e. share the same row, column, or diagonal). Backtracking explores partial assignments and backtracks upon conflict.

\section{Python Source Code (\texttt{Eight_Queen.py})}
\begin{lstlisting}[language=Python]
def is_safe(board, row, col):
    #Check current column in previous rows
    for i in range(row):
        if board[i][col] == 'Q':
            return False
    
    #Checkk upper-left diagonal
    i,j=row-1,col-1
    while i>=0 and j>=0:
        if board[i][j] == 'Q':
            return False
        i -= 1
        j -= 1
        
    #Check upper-right diagonal
    i,j = row-1,col+1
    while i>=0 and j<8:
        if board[i][j] == 'Q':
            return False
        i -= 1
        j += 1
    return True

def solve_8queens(board, row):
    #Base case: All 8 queens are placed
    if row == 8:
        return True
    
    for col in range(8):
        if is_safe(board, row, col):
            board[row][col] = 'Q'
            if solve_8queens(board, row+1):
                return True
            board[row][col] = '.'
            
    return False

#Initialize 8x8 board
board = [['.' for _ in range(8)] for _ in range(8)]

if solve_8queens(board,0):
    print("Solution for 8-Queens Problem:\n")
    for row in board:
        print(" ".join(row))
        
else:
    print("No Solution Exists.")
\end{lstlisting}

\section{Input Specifications}
None (Autonomous backtracking search placement on an 8x8 chessboard)

\section{Execution Output}
\begin{tcolorbox}[colback=outputback,colframe=outputframe,title={\textbf{Terminal Output - 8-Queens Problem}}]
\begin{verbatim}
Solution for 8-Queens Problem:

Q . . . . . . .
. . . . Q . . .
. . . . . . . Q
. . . . . Q . .
. . Q . . . . .
. . . . . . Q .
. Q . . . . . .
. . . Q . . . .
\end{verbatim}
\end{tcolorbox}

\section{Result and Conclusion}
The program for \textbf{8-Queens Problem} was implemented, tested, and validated successfully in Python. The outputs obtained comply with theoretical state-space traversal and algorithm complexity expectations.

\newpage

\chapter{Experiment 9: Water Jug Problem}

\section{Aim / Objective}
To solve the Water Jug problem using State-Space Search (BFS) to measure a target quantity of water using two jugs.

\section{Algorithm and Theory}
The Water Jug Problem is a state space search problem where states are represented as ordered pairs $(j_1, j_2)$ of water contents. Valid production rules include filling, emptying, and pouring between jugs until the target measurement is achieved.

\section{Python Source Code (\texttt{Water_jug.py})}
\begin{lstlisting}[language=Python]
from collections import deque

def water_jug(cap1, cap2, target_state):
    visited = set()
    queue = deque()  # Start with both jugs empty
    queue.append((0, 0, []))  # (jug1, jug2, path)
    while queue:
        j1, j2, path = queue.popleft()  

        if (j1, j2) in visited:
            continue
        visited.add((j1, j2))

        current_path = path + [(j1, j2)]

        if (j1, j2) == target_state:
            return current_path

        next_moves = [
            (cap1, j2),  # Fill jug1
            (j1, cap2),  # Fill jug2
            (0, j2),     # Empty jug1
            (j1, 0),     # Empty jug2
            (j1 - min(j1, cap2 - j2), j2 + min(j1, cap2 - j2)),  # Pour jug1 to jug2
            (j1 + min(j2, cap1 - j1), j2 - min(j2, cap1 - j1))   # Pour jug2 to jug1
        ]
        for move in next_moves:
            if move not in visited:
                queue.append((move[0], move[1], current_path))

    return None  # No solution found

target_goal = (4,2)
solution = water_jug(4, 3, target_goal)

if solution:
    print(f"steps to reach the target state {target_goal}:")
    for step in solution:
        print(step)
else:
    print("No solution exists.")
\end{lstlisting}

\section{Input Specifications}
None (Jug 1 Capacity: 4L, Jug 2 Capacity: 3L, Target state: (4, 2))

\section{Execution Output}
\begin{tcolorbox}[colback=outputback,colframe=outputframe,title={\textbf{Terminal Output - Water Jug Problem}}]
\begin{verbatim}
steps to reach the target state (4, 2):
(0, 0)
(0, 3)
(3, 0)
(3, 3)
(4, 2)
\end{verbatim}
\end{tcolorbox}

\section{Result and Conclusion}
The program for \textbf{Water Jug Problem} was implemented, tested, and validated successfully in Python. The outputs obtained comply with theoretical state-space traversal and algorithm complexity expectations.

\newpage

\chapter{Experiment 10: Tic-Tac-Toe Game}

\section{Aim / Objective}
To implement a two-player interactive Tic-Tac-Toe game board with state representation and win/draw validation.

\section{Algorithm and Theory}
Tic-Tac-Toe is a 2-player deterministic, perfect-information zero-sum game played on a $3 \times 3$ grid. The program maintains board state transitions and tests terminal win conditions along all rows, columns, and diagonals.

\section{Python Source Code (\texttt{Tic_Tac_Toe.py})}
\begin{lstlisting}[language=Python]
def print_board(board):
    for row in board:
        print(" | ".join(row))
        print("-" * 9)


def check_winner(board):
    # Check rows and columns
    for i in range(3):
        if board[i][0] == board[i][1] == board[i][2] != " ":
            return True

        if board[0][i] == board[1][i] == board[2][i] != " ":
            return True

    # Check diagonals
    if board[0][0] == board[1][1] == board[2][2] != " ":
        return True

    if board[0][2] == board[1][1] == board[2][0] != " ":
        return True

    return False


board = [[" " for _ in range(3)] for _ in range(3)]
current_player = "X"

for turn in range(9):
    print(f"\nPlayer {current_player}'s turn:")
    print_board(board)

    row = int(input("Enter row (0-2): "))
    col = int(input("Enter column (0-2): "))

    if board[row][col] == " ":
        board[row][col] = current_player

        if check_winner(board):
            print_board(board)
            print(f"\nPlayer {current_player} wins!")
            break

        current_player = "O" if current_player == "X" else "X"

    else:
        print("Cell already occupied. Try again.")

else:
    print_board(board)
    print("\nGame Draw!")
\end{lstlisting}

\section{Input Specifications}
Sequential row and column moves: (0,0), (1,0), (0,1), (1,1), (0,2) yielding a win for Player X

\section{Execution Output}
\begin{tcolorbox}[colback=outputback,colframe=outputframe,title={\textbf{Terminal Output - Tic-Tac-Toe Game}}]
\begin{verbatim}
Player X's turn:
  |   |  
---------
  |   |  
---------
  |   |  
---------
Enter row (0-2): Enter column (0-2): 
Player O's turn:
X |   |  
---------
  |   |  
---------
  |   |  
---------
Enter row (0-2): Enter column (0-2): 
Player X's turn:
X |   |  
---------
O |   |  
---------
  |   |  
---------
Enter row (0-2): Enter column (0-2): 
Player O's turn:
X | X |  
---------
O |   |  
---------
  |   |  
---------
Enter row (0-2): Enter column (0-2): 
Player X's turn:
X | X |  
---------
O | O |  
---------
  |   |  
---------
Enter row (0-2): Enter column (0-2): X | X | X
---------
O | O |  
---------
  |   |  
---------

Player X wins!
\end{verbatim}
\end{tcolorbox}

\section{Result and Conclusion}
The program for \textbf{Tic-Tac-Toe Game} was implemented, tested, and validated successfully in Python. The outputs obtained comply with theoretical state-space traversal and algorithm complexity expectations.

\newpage

\chapter{Experiment 11: Minimax Algorithm}

\section{Aim / Objective}
To implement the Minimax decision rule for adversarial game tree evaluation.

\section{Algorithm and Theory}
The Minimax algorithm is a recursive decision strategy in game theory and artificial intelligence. The maximizing player seeks to maximize the minimum payoff, while the minimizing player minimizes the opponent's maximum gain.

\section{Python Source Code (\texttt{minmax.py})}
\begin{lstlisting}[language=Python]
def minmax(depth,nodeIndex, isMax, scores, height):
    if depth == height:
        return scores[nodeIndex]
    
    if isMax:
        return max(minmax(depth + 1, nodeIndex * 2, False, scores, height),
                   minmax(depth + 1, nodeIndex * 2 + 1, False, scores, height))
    else:
        return min(minmax(depth + 1, nodeIndex * 2, True, scores, height),
                   minmax(depth + 1, nodeIndex * 2 + 1, True, scores, height))

#main program
scores = list(map(int, input("Enter 8 leaf node values: ").split()))
height = 3
result = minmax(0, 0, True, scores, height)
print("The optimal value is:", result)
\end{lstlisting}

\section{Input Specifications}
8 leaf node values: 3 5 6 9 1 2 0 -1 at tree height 3

\section{Execution Output}
\begin{tcolorbox}[colback=outputback,colframe=outputframe,title={\textbf{Terminal Output - Minimax Algorithm}}]
\begin{verbatim}
Enter 8 leaf node values: The optimal value is: 5
\end{verbatim}
\end{tcolorbox}

\section{Result and Conclusion}
The program for \textbf{Minimax Algorithm} was implemented, tested, and validated successfully in Python. The outputs obtained comply with theoretical state-space traversal and algorithm complexity expectations.

\newpage

\chapter{Experiment 12: Alpha-Beta Pruning}

\section{Aim / Objective}
To implement Alpha-Beta Pruning to prune non-promising branches from the Minimax game tree and reduce search complexity.

\section{Algorithm and Theory}
Alpha-Beta pruning optimizes Minimax search without altering the final optimal decision. It maintains two values: $\alpha$ (the highest score guaranteed to the maximizer) and $\beta$ (the lowest score guaranteed to the minimizer). Pruning occurs when $\beta \le \alpha$.

\section{Python Source Code (\texttt{Alpha_Beta_pruning.py})}
\begin{lstlisting}[language=Python]
import math
def alpha_beta(depth, nodeIndex, maximizingPlayer, values, alpha, beta, height):
    if depth==height:
        return values[nodeIndex]

    if maximizingPlayer:
        best = -math.inf
        for i in range(2):
           value = alpha_beta(depth+1, nodeIndex*2+i, False, values, alpha, beta, height)
           best = max(best,value)
           alpha = max(alpha, best)
           
           if beta <= alpha:
               break
        return best
    else:
        best = math.inf
        for i in range(2):
            value = alpha_beta(depth+1, nodeIndex*2+i, True, values, alpha, beta, height)
            best = min(best,value)
            beta = min(beta, best)
            
            if beta <= alpha:
                break
        return best

#Main Program
values = list(map(int, input("Enter 8 leaf node values: ").strip().split()))
height = 3
result = alpha_beta(0, 0, True, values, -math.inf, math.inf, height)
print("\nOptimal value:", result)
\end{lstlisting}

\section{Input Specifications}
8 leaf node values: 3 5 6 9 1 2 0 -1 with alpha=-infinity, beta=+infinity

\section{Execution Output}
\begin{tcolorbox}[colback=outputback,colframe=outputframe,title={\textbf{Terminal Output - Alpha-Beta Pruning}}]
\begin{verbatim}
Enter 8 leaf node values: 
Optimal value: 5
\end{verbatim}
\end{tcolorbox}

\section{Result and Conclusion}
The program for \textbf{Alpha-Beta Pruning} was implemented, tested, and validated successfully in Python. The outputs obtained comply with theoretical state-space traversal and algorithm complexity expectations.

\newpage

\chapter{Experiment 13: Blocks World Problem}

\section{Aim / Objective}
To solve the Blocks World planning problem using state-space BFS search to transition from initial block stacks to goal configuration.

\section{Algorithm and Theory}
The Blocks World problem is a benchmark planning domain in Artificial Intelligence consisting of planar blocks placed on a table or stacked on top of one another. The objective is to produce a sequence of valid unstack/stack moves to transform an initial state into a desired goal state.

\section{Python Source Code (\texttt{block_world.py})}
\begin{lstlisting}[language=Python]
from collections import deque


def is_clear(state, block):
    for b, position in state:
        if position == block:
            return False
    return True


def generate_moves(state, blocks):
    moves = []

    for block in blocks:
        if not is_clear(state, block):
            continue

        current_position = None
        for b, position in state:
            if b == block:
                current_position = position
                break

        # Move block to table
        if current_position != "Table":
            new_state = set(state)
            new_state.remove((block, current_position))
            new_state.add((block, "Table"))

            moves.append(
                (
                    frozenset(new_state),
                    f"Move {block} from {current_position} to Table"
                )
            )

        # Move block onto another clear block
        for destination in blocks:
            if block == destination:
                continue

            if current_position == destination:
                continue

            if not is_clear(state, destination):
                continue

            new_state = set(state)
            new_state.remove((block, current_position))
            new_state.add((block, destination))

            moves.append(
                (
                    frozenset(new_state),
                    f"Move {block} from {current_position} to {destination}"
                )
            )

    return moves


def blocks_world(initial_state, goal_state):
    queue = deque()
    queue.append((frozenset(initial_state), []))

    visited = set()
    visited.add(frozenset(initial_state))

    while queue:
        current_state, path = queue.popleft()

        if current_state == frozenset(goal_state):
            return path

        blocks = [block for block, _ in current_state]
        moves = generate_moves(current_state, blocks)

        for new_state, move in moves:
            if new_state not in visited:
                visited.add(new_state)
                queue.append((new_state, path + [move]))

    return None


# ---------------- MAIN PROGRAM ----------------

print("===== BLOCK WORLD PROBLEM =====")

# Number of blocks
n = int(input("Enter number of blocks: "))

blocks = []

print("\nEnter block names:")
for i in range(n):
    block = input(f"Block {i + 1}: ")
    blocks.append(block)

# Initial state
print("\n===== INITIAL STATE =====")
print("Enter the position of each block.")
print("Enter 'Table' if the block is on the table.")

initial_state = set()

for block in blocks:
    position = input(f"Where is {block} initially? ")
    initial_state.add((block, position))

# Goal state
print("\n===== GOAL STATE =====")
print("Enter the desired position of each block.")

goal_state = set()

for block in blocks:
    position = input(f"Where should {block} be in the goal state? ")
    goal_state.add((block, position))

# Display states
print("\nInitial State:")
for block, position in sorted(initial_state):
    print(block, "->", position)

print("\nGoal State:")
for block, position in sorted(goal_state):
    print(block, "->", position)

# Solve
solution = blocks_world(initial_state, goal_state)

# Display solution
if solution:
    print("\n===== SOLUTION =====")

    for i, move in enumerate(solution, 1):
        print(f"Step {i}: {move}")

    print("\nTotal moves:", len(solution))
else:
    print("\nNo solution found.")
\end{lstlisting}

\section{Input Specifications}
3 blocks (A, B, C); Initial: A on Table, B on A, C on Table; Goal: A on Table, B on Table, C on B

\section{Execution Output}
\begin{tcolorbox}[colback=outputback,colframe=outputframe,title={\textbf{Terminal Output - Blocks World Problem}}]
\begin{verbatim}
===== BLOCK WORLD PROBLEM =====
Enter number of blocks: 
Enter block names:
Block 1: Block 2: Block 3: 
===== INITIAL STATE =====
Enter the position of each block.
Enter 'Table' if the block is on the table.
Where is A initially? Where is B initially? Where is C initially? 
===== GOAL STATE =====
Enter the desired position of each block.
Where should A be in the goal state? Where should B be in the goal state? Where should C be in the goal state? 
Initial State:
A -> Table
B -> A
C -> Table

Goal State:
A -> Table
B -> Table
C -> B

===== SOLUTION =====
Step 1: Move B from A to Table
Step 2: Move C from Table to B

Total moves: 2
\end{verbatim}
\end{tcolorbox}

\section{Result and Conclusion}
The program for \textbf{Blocks World Problem} was implemented, tested, and validated successfully in Python. The outputs obtained comply with theoretical state-space traversal and algorithm complexity expectations.

\newpage

\end{document}
